
\chapter{Wstęp}

W ciągu ostatniej dekady znacząco wzrosło zainteresowanie tematem projektowania płytek PCB, a rozwój narzędzi oraz
bibliotek elementów sprawił, że znacząco obniżył się próg wejścia do tej dziedziny. W 2016 roku platforma EasyEDA
chwaliła się na swojej stronie internetowej liczbą 50 tysięcy projektów. W 2026 roku nadal jest to 50, ale nie
tysięcy, lecz milionów projektów. Widać zatem, że zapotrzebowanie na unikalne płytki PCB stale rosło. Wspomniane
obniżenie progu wejścia skutkuje tym, że do zaprojektowania płytki nie trzeba już być doświadczonym inżynierem,
a wystarczy znać podstawy elektrotechniki oraz wytyczne projektowania płytek. Tu do wyznaczenia trendu posłuży mi
liczba nowych projektów na GitHubie stworzonych na platformie KiCad. W 2012 roku były to zaledwie 62 projekty,
w latach 2020--2022 liczba ta trzymała się na poziomie 1500, a w roku 2023 było to prawie 2000 nowych projektów.

Rozwój technologii oraz większa świadomość użytkowników nie oznaczają jednak braku nowych wyzwań. Coraz bardziej
skomplikowane urządzenia oznaczają coraz bardziej złożone projekty, a w erze miniaturyzacji dodatkowym wyzwaniem
jest rozsądne rozmieszczenie elementów na płytce. To pociąga za sobą kolejne konsekwencje. Podczas gdy dla prostych
projektów przestrzeganie zasad projektowania często można pominąć, a ewentualne dostosowanie się do nich nie wymaga
dużego nakładu czasowego, tak w przypadku większych projektów odgrywa ono jedną z kluczowych ról w prawidłowym
funkcjonowaniu płytki i często oznacza godziny pracy fachowca. Co prawda platforma EasyEDA przygotowała dla
użytkowników narzędzie do automatycznego rozmieszczania elementów na płytce, podziału na warstwy i wyznaczania
ścieżek, jednak w praktyce, ze względu na jakość rozwiązań, nie korzysta się z niego.

Celem pracy jest przeprowadzenie badań nad skutecznością rozwiązań sztucznej inteligencji w tym właśnie zakresie.

Struktura pracy jest następująca:
\begin{itemize}
    \item w rozdziale 2 przedstawiono podstawowe zagadnienia oraz aktualny stan wiedzy w oparciu o przegląd literatury;
    \item rozdział 3 zawiera informacje na temat obiektu oraz metodologii badania;
    \item rozdział 4 poświęcony jest samym badaniom;
    \item rozdział 5 przedstawia wyniki badań;
    \item w rozdziale 6 sformułowane zostały wnioski dotyczące badań;
    \item rozdział 7 stanowi podsumowanie pracy.
\end{itemize}
